% !TEX root = ../ms-thesis.tex

\begin{acknowledgements}
衷心感谢邬东璠老师在选题、数据方法与论文写作全过程中的悉心指导。从开题到中期，再到本轮论文修订，老师一次次地将研究方向拉回``问题---潜力---对策''的主线，既避免了研究在空间生产理论等高阶议题上的过度拔高，也避免了研究停留在``当地人也知道''的浅层描述上。对桥梁选择、数据口径、错位指数、九江专题等具体环节的意见，几乎构成了本稿每一个章节的骨架。

感谢家人与朋友一直以来的理解与支持，使得在本科毕设的这段时光里，能够相对专注地完成这项研究。
\end{acknowledgements}
